\chapter{Prefazione}

\introduction{Andrea Calì}{
    \textit{L'obiettivo del corso è mostrare come due programmi, ovunque si trovino nel mondo,
    riescano a scambiarsi informazione. Si parte dalle applicazioni che usiamo ogni giorno e si
    scende, uno strato alla volta, fino ai bit che viaggiano sul filo: una rete non è una scatola
    magica, ma una pila di protocolli, ciascuno con un problema preciso da risolvere. Lungo il
    percorso si scrivono programmi in C sui socket e si chiude con la sicurezza e la privatezza.}
  }
\pagestyle{plain}

% ══════════════════════════════════════════════════════════════
\section{Come è organizzata questa dispensa}
% ══════════════════════════════════════════════════════════════

Il corso segue un approccio \textbf{top-down}: si parte dal livello che l'utente percepisce
(le applicazioni: il Web, la posta, il DNS) e si scende progressivamente verso l'hardware
(il livello di collegamento e il livello fisico). La dispensa rispetta lo stesso ordine ed è
divisa in sette parti.

\begin{center}
\begin{tikzpicture}[every node/.style={font=\small\sffamily}]
  \node[box, minimum width=2.6cm, minimum height=2.4cm, fill=lphy!40] (f) at (0,0) {Parte I\\[2pt]Fondamenti\\[2pt]\scriptsize(modello\\\scriptsize a strati)};
  \node[lay app, minimum width=4.4cm] (a) at (4.6,0.93) {Parte II --- Applicazione};
  \node[lay tra, minimum width=4.4cm] (t) at (4.6,0.31) {Parte III --- Trasporto};
  \node[lay net, minimum width=4.4cm] (n) at (4.6,-0.31) {Parte IV --- Rete};
  \node[lay link, minimum width=4.4cm] (l) at (4.6,-0.93) {Parte V --- Collegamento};
  \node[box, minimum width=2.6cm, minimum height=2.4cm, fill=netred!10] (s) at (9.2,0) {Parte VI\\[2pt]Sicurezza};
  \node[box, minimum width=2.6cm, minimum height=2.4cm, fill=netgreen!10] (e) at (12.4,0) {Parte VII\\[2pt]Esercizi};
  \draw[-{Stealth[length=3mm]}, very thick, primary] (f.east) -- (a.west |- f.east);
  \draw[-{Stealth[length=3mm]}, very thick, primary] (l.east |- s.west) -- (s.west);
  \draw[-{Stealth[length=3mm]}, very thick, primary] (s.east) -- (e.west);
  \draw[-{Stealth[length=2.5mm]}, thick, primary] ([xshift=0.25cm]a.north east) -- node[right, lbl, rotate=-90, anchor=north] {dall'alto in basso} ([xshift=0.25cm]l.south east);
\end{tikzpicture}
\end{center}

\begin{itemize}
    \item \textbf{Parte I --- Fondamenti}: che cos'è una rete, com'è fatta, come si classifica, e
          soprattutto il \emph{modello a strati} che regge tutto il resto del corso.
    \item \textbf{Parte II --- Livello applicativo}: architetture client-server e P2P, HTTP,
          posta elettronica, DNS e la programmazione di rete in C con i socket.
    \item \textbf{Parte III --- Livello di trasporto}: UDP, il problema del trasferimento affidabile,
          TCP in tutte le sue componenti (numeri di sequenza, timeout, connessione, controllo di
          flusso e di congestione).
    \item \textbf{Parte IV --- Livello di rete}: router, indirizzamento IP, subnetting e CIDR,
          DHCP, NAT, IPv6, ICMP e gli algoritmi di routing.
    \item \textbf{Parte V --- Livello di collegamento}: framing, rilevazione degli errori (parità,
          CRC, distanza di Hamming), protocolli ad accesso multiplo, Ethernet, switch e ARP.
    \item \textbf{Parte VI --- Sicurezza}: attacchi, crittografia simmetrica e asimmetrica,
          integrità dei messaggi, autenticazione, SSL/TLS, firewall e IDS.
    \item \textbf{Parte VII --- Esercizi}: esercizi svolti divisi per argomento, esercitazioni di
          laboratorio con Wireshark e domande di ripasso.
\end{itemize}

\info{Una dispensa che cresce con il corso}{
    La dispensa viene aggiornata dopo ogni lezione del corso di \emph{Reti e Programmazione
    Distribuita}. La Parte I (lezioni 1 e 2) e il primo capitolo della Parte II con la sezione sul P2P
    (lezione 3) seguono già le trasparenze, le registrazioni e le lavagne del nuovo corso; i capitoli
    successivi sono basati sul materiale dell'edizione precedente (\emph{Reti di Calcolatori I})
    e verranno rivisti man mano che le lezioni procedono.
}

% ══════════════════════════════════════════════════════════════
\section{Convenzioni tipografiche}
% ══════════════════════════════════════════════════════════════

Nel testo si usano alcuni riquadri ricorrenti:

\dfn{Definizione}{
    I riquadri rossi contengono le \textbf{definizioni} formali: sono i concetti da saper enunciare
    con precisione all'esame.
}

\info{Nota}{
    I riquadri verdi contengono \textbf{approfondimenti}, analogie e osservazioni utili a capire
    meglio, ma non strettamente indispensabili.
}

\warning{Attenzione}{
    I riquadri gialli segnalano \textbf{errori tipici} o punti in cui è facile confondersi.
}

\begin{esempio}{Struttura di un esempio}
Gli esempi sono racchiusi tra due righe orizzontali sottili e sono numerati all'interno di ogni capitolo.
\end{esempio}

\begin{esercizio}{Struttura di un esercizio}
Gli esercizi sono racchiusi tra righe doppie. Nella Parte VII si trovano tutti gli esercizi svolti.
\end{esercizio}

Nelle figure i livelli della pila protocollare hanno sempre lo stesso colore, così da
riconoscerli a colpo d'occhio:

\begin{center}
\begin{tikzpicture}[node distance=0.15cm]
  \node[lay app, minimum width=2.3cm] (a) {Applicazione};
  \node[lay tra, minimum width=2.3cm, right=of a] (t) {Trasporto};
  \node[lay net, minimum width=2.3cm, right=of t] (n) {Rete};
  \node[lay link, minimum width=2.3cm, right=of n] (l) {Collegamento};
  \node[lay phy, minimum width=2.3cm, right=of l] (p) {Fisico};
\end{tikzpicture}
\end{center}

% ══════════════════════════════════════════════════════════════
\section{Testi di riferimento e materiale}
% ══════════════════════════════════════════════════════════════

Il corso fa riferimento principalmente a due testi, entrambi citati nelle lezioni:

\begin{itemize}
    \item \textbf{J.F. Kurose, K.W. Ross}, \emph{Computer Networking: A Top-Down Approach}. È il testo
          che dà l'impostazione top-down del corso e da cui provengono molti esempi su HTTP, TCP e IP.
    \item \textbf{A.S. Tanenbaum, N. Feamster, D.J. Wetherall}, \emph{Computer Networks}. Più prolisso
          ma ricchissimo di esempi, anche non strettamente tecnici; usato soprattutto per la
          classificazione delle reti, il livello di collegamento e la crittografia.
\end{itemize}

Le trasparenze sono in inglese (la terminologia di rete è nata in inglese ed è più semplice
impararla così); in questa dispensa i termini tecnici sono riportati sia in italiano sia in
inglese. Le trasparenze sono un \textbf{supporto} alla lezione, non un sostituto dei libri.

\info{Il laboratorio}{
    Per le esercitazioni viene fornito un \textbf{container Docker} già pronto, con tutto il software
    necessario (compilatore, strumenti di rete, Wireshark). Un container è una sorta di macchina
    virtuale leggera che condivide il kernel con la macchina ospite: se qualcosa si rompe, basta
    buttarlo via e ricrearlo, senza toccare il proprio sistema.
}
